\chapter{Exchange Rates}\label{ch:b3:exchange-rates}

An exchange rate is the price of one money in terms of another, and it is set in the
largest and most liquid market in the world. Its theory rests on two arbitrage
conditions --- one in goods, one in assets --- and the evidence finds one of them
holding slowly and the other not at all. This chapter tests \omterm{def:b3:exchange-rates:ppp}{purchasing-power parity}
over a century and a half, explains why rich countries are dear, shows that
\omterm{def:b3:exchange-rates:interest-parity}{covered interest parity} holds while uncovered parity fails, derives the
exchange rate in the monetary model and its \omterm{def:b3:exchange-rates:overshooting}{overshooting}, and ends with the
result that a random walk forecasts exchange rates better than the models.

\begin{openingchart}{Price level index (GDP; United States = 100) against GDP per
person (constant 2021 PPP dollars, log scale), 197 countries, 2023, with the fitted
line of log price level on log income.}{World Bank, World Development Indicators
(PA.NUS.GDP.PLI, NY.GDP.PCAP.PP.KD); our computation}
\begin{tikzpicture}
\begin{axis}[omscatter, xmode=log, log ticks with fixed point, width=0.66\linewidth,
  height=0.44\linewidth, xmin=800, xmax=150000, ymin=0, ymax=130,
  xlabel={GDP per person}, ylabel={Price level (United States = 100)}]
  \addplot table[x=y, y=pli] {figdata/bachelor-3/13-exchange-rates/out/pli.csv};
  \addplot[omThm, thick, no marks, smooth] table[x=y, y=pli] {figdata/bachelor-3/13-exchange-rates/out/pli_fit.csv};
\end{axis}
\end{tikzpicture}
\end{openingchart}

% figdata: pli_stats
At market exchange rates, the same basket of final goods cost 111 in Switzerland
when it cost 100 in the United States, 51 in China, 25 in India and 19 in Egypt.
Across countries, a country twice as rich has a price level about 15~\% higher
(the slope is 0.20, $2^{0.20} = 1.15$), and income alone accounts for a third of the
variation. The law of one price does not make price levels equal.

\begin{recall}
The Year 1 volume defined the nominal exchange rate $S$ (here, home currency per
unit of foreign currency, so that a rise of $S$ is a depreciation of the home
currency) and the real exchange rate $Q = SP^*/P$. The Year 2 volume proved the law of
one price for assets and modelled expectations.
\end{recall}

\begin{notation}[Symbols of this chapter]\label{not:b3:exchange-rates:symbols}
Lower-case letters are logarithms: $s = \ln S$, $p = \ln P$, $q = s + p^* - p$, $m$ money.
$F$ is the one-year forward rate, $i$ and $i^*$ the home and foreign one-year interest
rates, $\lambda$ the interest semi-elasticity of money demand and $\theta$ the speed at
which a deviation closes. A star marks the foreign country.
\end{notation}

\section{Purchasing-power parity}\label{sec:b3:exchange-rates:ppp}

\begin{definition}[Absolute and relative purchasing-power parity]\label{def:b3:exchange-rates:ppp}
\index{purchasing-power parity}
\emph{Absolute purchasing-power parity} holds when a common basket costs the same in
two countries once converted at the exchange rate: $S = P/P^*$, a real exchange rate
of one. \emph{Relative purchasing-power parity} holds when the exchange rate moves
with the inflation differential: $\Delta s = \pi - \pi^*$, so that the real exchange rate is
constant.
\end{definition}

Absolute parity implies relative parity (take log differences of $S = P/P^*$), but
not conversely: a constant real exchange rate different from one satisfies the
second only. Neither holds year by year. The question is whether deviations die
out.

\begin{definition}[Half-life of a deviation]\label{def:b3:exchange-rates:half-life}
\index{half-life of a deviation}
If the real exchange rate follows $q_t - \bar q = \rho(q_{t-1} - \bar q) + \varepsilon_t$ with $0 < \rho < 1$, the
\emph{half-life} of a deviation is the time it takes for half of it to disappear,
$\ln 0.5/\ln \rho$ periods.
\end{definition}

\begin{omchart}{Real exchange rates against the dollar, 1870--2020 (log deviation
from each country's 1870--2020 mean, percent; a rise is a real depreciation against
the dollar). The war years 1914--1919 and 1939--1949, of controlled official rates
and wartime inflation, are left out.}{Jordà, Schularick and Taylor, Macrohistory Database (xrusd, cpi); our
computation}
\begin{tikzpicture}
\begin{axis}[omtime, width=0.8\linewidth, height=0.46\linewidth, xmin=1865, xmax=2025,
  ymin=-130, ymax=100, ylabel={Percent}, unbounded coords=jump, legend columns=5,
  legend style={at={(0.5,1.03)}, anchor=south, font=\footnotesize}]
  \addplot[omDef, thick, mark=none] table[x=year, y=GBR] {figdata/bachelor-3/13-exchange-rates/out/real_rates.csv};
  \addplot[omThm, thick, mark=none] table[x=year, y=JPN] {figdata/bachelor-3/13-exchange-rates/out/real_rates.csv};
  \addplot[black, thick, mark=none] table[x=year, y=CHE] {figdata/bachelor-3/13-exchange-rates/out/real_rates.csv};
  \addplot[omDef, thick, dashed, mark=none] table[x=year, y=CAN] {figdata/bachelor-3/13-exchange-rates/out/real_rates.csv};
  \addplot[gray, thick, dotted, mark=none] table[x=year, y=AUS] {figdata/bachelor-3/13-exchange-rates/out/real_rates.csv};
  \legend{United Kingdom, Japan, Switzerland, Canada, Australia}
\end{axis}
\end{tikzpicture}
\end{omchart}

% figdata: half_life_stats
The British, Canadian and Australian real rates wander around a stable level.
The Japanese and Swiss ones do not: both currencies appreciated in real terms by
more than 100 log points over the period. Estimating $\rho$ for each of the 17
economies of the database gives a median half-life of 5.5 years over 1870--2020 and
2.5 years over 1974--2020 (5.4 years over 1870--2020 without the war years), against a consensus of three to five years in the
studies surveyed in 1996 --- too slow to be explained by sticky prices alone.
% ledger: b3xr:rog.halflife
The short sample's shorter half-life should not be read as faster convergence:
least squares underestimates $\rho$ in short samples, and the bias shortens the
half-life.

\begin{model}{Purchasing-power parity}
\assumes Arbitrage in goods: a basket cheaper in one country is bought there and
sold in the other until prices are equal in a common currency.
\predicts Real exchange rates that revert to a constant level; nominal rates that
follow inflation differentials in the long run.
\tested Deviations die out, with half-lives of three to five years in long samples
(5.5 years for the median economy here); they are large in the short run, and
some real rates trend for a century. % ledger: b3xr:rog.halflife
\breaks \omterm{def:b3:ricardian-trade:iceberg}{Non-traded goods}, whose prices differ across countries; trade costs; the
short run, when nominal rates move far more than prices; trends in productivity.
\end{model}

\section{Productivity and price levels}\label{sec:b3:exchange-rates:balassa-samuelson}

\begin{definition}[Balassa--Samuelson effect]\label{def:b3:exchange-rates:balassa-samuelson}
\index{Balassa--Samuelson effect}
The \emph{Balassa--Samuelson effect} is the rise of a country's price level, relative
to others, that comes from faster productivity growth in its traded-goods sector
than in its non-traded sector.
\end{definition}

\begin{proposition}[Balassa--Samuelson]\label{prop:b3:exchange-rates:balassa-samuelson}
Let each country produce a traded good and a \omterm{def:b3:ricardian-trade:iceberg}{non-traded good} with labour alone,
$y_T = a_TL_T$ and $y_N = a_NL_N$; let labour move freely between sectors, the law of one
price hold for the traded good, and the price level be $P = P_T^{1-\gamma}P_N^\gamma$. Then
\[
\frac{P}{SP^*} = \left(\frac{a_T/a_N}{a_T^*/a_N^*}\right)^{\gamma}.
\]
\end{proposition}

\begin{proof}
With competitive firms, each sector pays labour its marginal product in value:
$W = P_Ta_T = P_Na_N$, so $P_N/P_T = a_T/a_N$. Then $P = P_T(P_N/P_T)^\gamma = P_T(a_T/a_N)^\gamma$, and the
same holds abroad. Dividing, with $P_T = SP_T^*$, gives the result.
\end{proof}

Productivity grows faster in manufacturing than in haircuts. A country whose
traded sector catches up raises the wage in both sectors, and the price of the
\omterm{def:b3:ricardian-trade:iceberg}{non-traded good} rises with it. Japan, the major economy whose income grew fastest
after the Second World War, is the canonical case of a sustained real appreciation
(the chart of real rates); the opening chart's slope is the cross-section.
% ledger: b3xr:rog.japan

\begin{measure}[How a price level index is made]
The price level index divides a country's purchasing-power-parity conversion
factor --- the number of local currency units that buy what one dollar buys in the
United States, estimated from the prices of comparable goods and services collected
in the International Comparison Program --- by its market exchange rate. An index
of 25 says that a dollar converted at the market rate buys four times as much there
as at home. The comparisons require that baskets be comparable, which is hardest
for services whose quality differs: a doctor's visit, a year of schooling.
\end{measure}

\section{Interest parity}\label{sec:b3:exchange-rates:interest-parity}

\begin{definition}[Covered and uncovered interest parity]\label{def:b3:exchange-rates:interest-parity}
\index{covered interest parity}\index{uncovered interest parity}
\emph{Covered interest parity} holds when $F/S = (1 + i)/(1 + i^*)$: a deposit at home earns
the same as a foreign deposit converted at the spot rate and converted back at a
forward rate agreed today. \emph{Uncovered interest parity} holds when the expected
depreciation equals the interest differential: $E_t s_{t+1} - s_t = i - i^*$.
\end{definition}

\begin{proposition}[Covered interest parity by arbitrage]\label{prop:b3:exchange-rates:cip}
If funds can be borrowed and lent at $i$ at home and $i^*$ abroad and forward
contracts traded without cost or default, then $F = S(1 + i)/(1 + i^*)$.
\end{proposition}

\begin{proof}
Borrow one unit of home currency at $i$, convert it to $1/S$ units of foreign
currency, lend them at $i^*$, and sell forward the $(1 + i^*)/S$ units due in a year at
$F$. Every price is known today; the strategy needs no capital, and at the end
yields $F(1 + i^*)/S - (1 + i)$ for certain. If this were positive, everyone would do it;
if negative, the reverse. No arbitrage requires it to be zero.
\end{proof}

Covered parity is the textbook arbitrage, yet in major currencies deviations from
it have been large, persistent and systematic, not explained by credit risk or
transaction costs, and strongest at quarter-ends, when banks' balance sheets are
reported: arbitrage needs balance sheet, and bank regulation made it costly.
% ledger: b3xr:dtv.cip
Uncovered parity is not an arbitrage: the investor who does not cover bears the
exchange risk. It is a prediction about expectations, and it can be tested. If
investors are risk-neutral and expectations rational, the regression of the
realised depreciation on the interest differential has a slope of one.

\begin{omchart}{Depreciation against the dollar over the next year (percent) against
the short-term interest differential with the United States (points), 17
economies, 1974--2019, with the fitted line and the uncovered-parity
line.}{Jordà, Schularick and Taylor, Macrohistory Database (xrusd, stir); our
computation}
\begin{tikzpicture}
\begin{axis}[omscatter, width=0.6\linewidth, height=0.44\linewidth, xmin=-14, xmax=14,
  ymin=-30, ymax=40, xlabel={Interest differential (points)},
  ylabel={Depreciation next year (\%)}, legend pos=north west,
  legend style={font=\footnotesize}]
  \addplot+[mark size=0.8pt] table[x=diff, y=dep] {figdata/bachelor-3/13-exchange-rates/out/uip.csv};
  \addplot[omThm, thick, no marks, sharp plot] table[x=diff, y=dep] {figdata/bachelor-3/13-exchange-rates/out/uip_fit.csv};
  \addplot[black, thick, dashed, no marks, sharp plot] table[x=diff, y=uip] {figdata/bachelor-3/13-exchange-rates/out/uip_fit.csv};
  \legend{Country-year, Fitted, Uncovered parity}
\end{axis}
\end{tikzpicture}
\end{omchart}

% figdata: uip_stats
The pooled slope is $-0.10$ (standard error 0.13), over 782 country-years: if
anything, high-interest currencies appreciated. The interest differential
explains essentially none of the variation in exchange rates (an $R^2$ under 0.001).

\begin{definition}[Forward premium puzzle; carry trade]\label{def:b3:exchange-rates:forward-premium}
\index{forward premium puzzle}\index{carry trade}
The \emph{forward premium puzzle} is the finding that the slope of the
uncovered-parity regression is below one and often negative. A \emph{carry trade}
borrows in low-interest currencies and lends in high-interest ones, uncovered.
\end{definition}

\begin{record}[The carry trade]
Under uncovered parity the \omterm{def:b3:exchange-rates:forward-premium}{carry trade} earns nothing on average: the
high-interest currency depreciates by the differential. In the data it has earned
large average payoffs, largely unrelated to standard measures of risk. One
explanation is a rare event: losses, or simply bad times, that come in states of
the world when investors value money most, and that a short sample may barely
contain.
% ledger: b3xr:bekr.carry
\end{record}

\section{The monetary model}\label{sec:b3:exchange-rates:monetary}

\begin{definition}[Monetary model of the exchange rate]\label{def:b3:exchange-rates:monetary-model}
\index{monetary model of the exchange rate}
The \emph{monetary model of the exchange rate} combines money-market equilibrium in
each country, $m - p = \phi y - \lambda i$, with \omterm{def:b3:exchange-rates:ppp}{purchasing-power parity} and \omterm{def:b3:exchange-rates:interest-parity}{uncovered interest
parity}.
\end{definition}

\begin{proposition}[The exchange rate as an asset price]\label{prop:b3:exchange-rates:monetary}
In the monetary model, with fundamentals $f_t = (m_t - m_t^*) - \phi(y_t - y_t^*)$ and no bubble,
\[
s_t = \frac{1}{1 + \lambda}\sum_{k=0}^{\infty}\left(\frac{\lambda}{1 + \lambda}\right)^{k}E_tf_{t+k}.
\]
\end{proposition}

\begin{proof}
Subtracting the foreign money market from the home one, $p - p^* = f + \lambda(i - i^*)$.
Parity in goods gives $s = p - p^*$, and in assets $i - i^* = E_ts_{t+1} - s_t$, so $s_t = f_t +
\lambda(E_ts_{t+1} - s_t)$, or $s_t = (f_t + \lambda E_ts_{t+1})/(1 + \lambda)$. Substituting the same equation
for $E_ts_{t+1}$, and so on $K$ times (a recurrence of the kind solved in the grades
10--12 volume), gives the sum up to $K$ plus $(\lambda/(1 + \lambda))^{K+1}E_ts_{t+K+1}$, which
vanishes when there is no bubble.
\end{proof}

The exchange rate is an asset price: it moves today with news about future
money and output. A rise in expected money growth depreciates the currency now,
before any money is printed (\cref{exo:b3:exchange-rates:11}).

\section{Overshooting}\label{sec:b3:exchange-rates:overshooting}

\begin{definition}[Overshooting]\label{def:b3:exchange-rates:overshooting}
\index{overshooting}\index{Dornbusch model}
A currency \emph{overshoots} when its response on impact to a permanent shock
exceeds its long-run response. In the \emph{Dornbusch model}, prices adjust slowly
while the exchange rate jumps, and a monetary expansion makes the currency
overshoot.
\end{definition}

\begin{proposition}[Overshooting]\label{prop:b3:exchange-rates:dornbusch}
Take the home money market $m - p = -\lambda i$ (output fixed), \omterm{def:b3:exchange-rates:interest-parity}{uncovered interest parity}
$i = i^* + E\dot s$, and prices that adjust slowly, so that the exchange rate is expected
to close its gap with the long-run rate $\bar s$ at speed $\theta$: $E\dot s = \theta(\bar s - s)$. After a
permanent rise $dm$ in money, $d\bar s = dm$, and on impact $ds = dm(1 + 1/(\lambda\theta))$.
\end{proposition}

\begin{proof}
In the long run $i = i^*$ and $p$ rises with $m$, so with parity in goods $\bar s$ rises by
$dm$. On impact $p$ has not moved, so the money market requires $i$ to fall by
$dm/\lambda$. Uncovered parity then requires an expected appreciation of $dm/\lambda$: $\theta(s -
\bar s) = dm/\lambda$, so $s - \bar s = dm/(\lambda\theta)$ and $ds = d\bar s + dm/(\lambda\theta)$.
\end{proof}

\begin{omfigure}
\begin{tikzpicture}[econ, x=0.62cm, y=0.12cm]
  \draw[econ axes] (-2,0) -- (12.5,0) node[right] {years};
  \draw[econ axes] (0,0) -- (0,23) node[above] {\% change};
  \draw[guide] (-2,10) node[above right, font=\scriptsize] {long run} -- (12,10);
  \draw[curve, omDef] (-2,0) -- (0,0) -- (0,20);
  \draw[curve, omDef, domain=0:12, samples=60] plot (\x, {10 + 10*exp(-0.35*\x)});
  \node[omDef, font=\scriptsize, anchor=south west] at (2.2,14.8) {exchange rate $s$};
  \draw[curve, omThm, domain=0:12, samples=60] plot (\x, {10 - 10*exp(-0.35*\x)});
  \node[omThm, font=\scriptsize, anchor=north west] at (2.4,5.2) {price level $p$};
  \draw[curve, omThm] (-2,0) -- (0,0);
  \foreach \t in {0,2,...,12} \draw (\t,0) -- (\t,-0.8) node[below, font=\scriptsize] {\t};
  \foreach \v in {10,20} \draw (0,\v) -- (-0.15,\v) node[left, font=\scriptsize] {\v};
\end{tikzpicture}
\par\smallskip{\small A permanent 10~\% rise in money at year 0, with $\lambda\theta = 1$ and $\theta = 0.35$:
the currency depreciates by 20~\% at once, then appreciates toward its long-run
depreciation of 10~\% as prices rise.}
\end{omfigure}

\begin{model}{The Dornbusch overshooting model}
\assumes Sticky goods prices, flexible asset prices, \omterm{def:b3:exchange-rates:interest-parity}{uncovered interest parity},
a stable demand for money.
\predicts A monetary expansion depreciates the currency at once by more than in the
long run, and the currency then appreciates while prices rise; exchange rates more
volatile than prices.
\tested Floating exchange rates are an order of magnitude more volatile than
relative prices, as predicted. % ledger: b3xr:rog.volatility
But after
identified U.S. monetary expansions, the dollar depreciates persistently, and
its largest fall comes with a delay, not on impact. % ledger: b3xr:ee.delayed
\breaks Uncovered parity, on which the jump depends, fails in the data (the chart
above); money demand is unstable.
\end{model}

\section{Forecasting}\label{sec:b3:exchange-rates:forecasting}

\begin{record}[Meese and Rogoff, 1983]
The exchange-rate models of the 1970s were tested out of sample on the first years
of floating. For the dollar against the pound, the mark and the yen, and for the
trade-weighted dollar, a random walk --- tomorrow's rate equals today's ---
forecast as well as any estimated model at horizons of one to twelve months. The
structural models did poorly even when given the actual future values of their
explanatory variables. % ledger: b3xr:mr.rw
Short-term movements of nominal exchange rates remained, in the words of a later
survey, notoriously difficult to explain even after the fact.
% ledger: b3xr:rog.volatility
\end{record}

The two halves of this chapter fit. If the exchange rate is an asset price that
moves with news about the future, its changes are mostly unpredictable, as asset
prices' are; and if deviations from parity die out with a half-life of years, a
model that knows where the rate is going has little to say about next month.

\section{Exercises}\label{sec:b3:exchange-rates:exercises}

\begin{exercise}[$\star$]\label{exo:b3:exchange-rates:1}
In the opening chart, which country has the highest price level, and which the
lowest? Which country is much cheaper than its income predicts? By how much does
the fitted line say the price level rises when income doubles?
\end{exercise}

\begin{exercise}[$\star$]\label{exo:b3:exchange-rates:2}
Home inflation is 6~\% a year and foreign inflation 2~\%. What does relative PPP
predict for the exchange rate, approximately and exactly?
\end{exercise}

\begin{exercise}[$\star$]\label{exo:b3:exchange-rates:3}
The spot rate is 1.10 dalens per kest, the one-year dalen rate 4~\% and the kest rate
2~\%. Find the one-year forward rate. Is the dalen at a forward premium or discount?
\end{exercise}

\begin{exercise}[$\star$]\label{exo:b3:exchange-rates:4}
What evidence would contradict \omterm{def:b3:exchange-rates:ppp}{purchasing-power parity} in the long run?
\end{exercise}

\begin{exercise}[$\star\star$]\label{exo:b3:exchange-rates:5}
Estimate the half-life of each JST economy's real exchange rate against the dollar,
over 1870--2020 and over 1974--2020. Which economies have the longest half-lives over
the long sample, and why?
\end{exercise}

\begin{exercise}[$\star\star$]\label{exo:b3:exchange-rates:6}
In the Balassa--Samuelson model with $\gamma = 0.5$, home's ratio $a_T/a_N$ is twice
the foreign one. How much higher is home's price level? If home's traded
productivity then grows by 3~\% a year faster than abroad, with $a_N$ unchanged
everywhere, how fast does its real exchange rate appreciate?
\end{exercise}

\begin{exercise}[$\star\star$]\label{exo:b3:exchange-rates:7}
In the \omterm{def:b3:exchange-rates:overshooting}{Dornbusch model} with $\lambda = 2$ and $\theta = 0.25$, money rises by 10~\%. Find the
depreciation on impact, after one year and in the long run, and the change in the
home interest rate on impact.
\end{exercise}

\begin{exercise}[$\star\star$]\label{exo:b3:exchange-rates:8}
Estimate the uncovered-parity regression country by country on the JST data,
1974--2019. In how many economies is the slope negative? In how many can a slope of
one be rejected at 5~\%?
\end{exercise}

\begin{exercise}[$\star\star$]\label{exo:b3:exchange-rates:9}
What forecasting record would contradict the Meese--Rogoff result?
\end{exercise}

\begin{exercise}[$\star\star\star$]\label{exo:b3:exchange-rates:10}
In the \omterm{def:b3:exchange-rates:overshooting}{Dornbusch model} of the chapter, show that the gap $s - \bar s$ after the shock
and the price gap $\bar p - p$ satisfy $s - \bar s = (\bar p - p)/(\lambda\theta)$ at every date, and that the
overshoot shrinks as prices become more flexible.
\end{exercise}

\begin{exercise}[$\star\star\star$]\label{exo:b3:exchange-rates:11}
In the monetary model with $\lambda = 4$, relative output constant and relative money
growing at 2~\% a year, show that $s_t = f_t + \lambda\mu$, where $\mu$ is money growth. The central
bank announces that money growth will be 5~\% a year from now on. By how much does
the currency depreciate on the announcement?
\end{exercise}

\begin{exercise}[$\star\star\star$]\label{exo:b3:exchange-rates:12}
An investor borrows yen at 0.5~\% and lends Australian dollars at 5~\% for a year,
uncovered. Write the return as the interest differential minus the depreciation.
Find it if the Australian dollar falls by 3~\%, and if it falls by 25~\%. What
probability of the crash makes the expected return zero, if the other outcome is a
2~\% rise?
\end{exercise}

\section{Problem: The Dalen Against the Kest}\label{sec:b3:exchange-rates:problem}

\begin{problem}[{Weekend problem --- The dalen against the kest}]\label{pb:b3:exchange-rates:dalen}
The rate $S$ (dalens per kest) and Brenland's CPI and policy rate come from its
accounts. Assume that Kestrian prices rise by 2.7~\% a year and that Kestria's
interest rate is 3.0~\%. The real rate $Q = SP_K/P_B$ is indexed to 2015 = 100:
% figdata: dalen_points
\begin{center}\small
\pgfplotstabletypeset[col sep=comma, columns={year,rate,cpi,cpi_k,real,policy_rate},
  columns/year/.style={column name={Year}, int detect, column type=l},
  columns/rate/.style={column name={$S$}, fixed, fixed zerofill, precision=3},
  columns/cpi/.style={column name={$P_B$}, fixed, fixed zerofill, precision=1},
  columns/cpi_k/.style={column name={$P_K$}, fixed, fixed zerofill, precision=1},
  columns/real/.style={column name={$Q$}, fixed, fixed zerofill, precision=1},
  columns/policy_rate/.style={column name={$i_B$ (\%)}, fixed, fixed zerofill, precision=2},
  every head row/.style={after row=\hline}, every last row/.style={after row=\hline}]
  {figdata/bachelor-3/13-exchange-rates/out/dalen_points.csv}
\end{center}
An AR(1) of $\ln Q$ on 2000--2025 gives $\rho = 0.948$. In 2021, $S = 1.120$, $P_B = 115.9$ and
$P_K = 117.3$; in 2022, $S = 1.148$, $P_B = 127.7$, $P_K = 120.5$ and $i_B = 5.42~\%$; in
2023, $S = 1.183$. Over 2016--2020 Brenland's policy rate ranged from 2.23~\% to
2.72~\%. Kestria's has been 3.0~\% throughout. % figdata: dalen_stats, dalen

\medskip\noindent\textbf{Part I --- parity in goods.}
\begin{enumerate}
\item What change in $S$ from 2000 to 2025 does relative PPP predict? What happened?
\item In which listed year was the dalen weakest in real terms? By how much was it
weaker than in 2015?
\item Show that the table's $Q$ for 2025 follows from $S$, $P_K$ and $P_B$.
\item Find the half-life of a real deviation implied by $\rho$.
\item Compare it with the JST medians of 5.5 and 2.5 years. Why is an estimate from
26 years imprecise?
\item Does the estimate show that PPP fails between Brenland and Kestria?
\end{enumerate}

\medskip\noindent\textbf{Part II --- parity in assets.}
\begin{enumerate}[resume]
\item Find the one-year forward rate in 2022 by covered parity.
\item What depreciation of the dalen did uncovered parity predict for 2022--2023?
What happened?
\item A Kestrian investor borrowed kests at 3.0~\% in 2022 and lent dalens at 5.42~\% for
a year, uncovered. Find the return in percent of the sum borrowed.
\item Over 2016--2020 Brenland's rate was below Kestria's. What did uncovered parity
predict, and what did the dalen do?
\item Compare the Brenland record with the JST regression slope of $-0.10$.
\item Name two reasons why uncovered parity can fail with rational investors.
\end{enumerate}

\medskip\noindent\textbf{Part III --- a Kestrian monetary expansion.} Kestria's money
supply rises permanently by 10~\%. In Kestria's money market $\lambda = 4$, and prices
close a deviation with a half-life of 2.5 years.
\begin{enumerate}[resume]
\item What happens to $S$ in the long run?
\item Find $\theta$, with $e^{-\theta \times 2.5} = 0.5$.
\item Find the change in $S$ on impact.
\item By how many points does the dalen overshoot?
\item Find the change in $S$ one year after the shock.
\item By how much does Kestria's interest rate fall on impact?
\item What would happen with fully flexible prices?
\end{enumerate}

\medskip\noindent\textbf{Part IV --- a forecast contest.} One-year-ahead forecasts of
$S$ for 2016--2025: a random walk; uncovered parity, $S_t(1 + i_B - i_K)$; and a PPP
rule that closes the real gap at the speed estimated with data up to the forecast
date and, like the models Meese and Rogoff tested, is given the actual inflation of
the year it forecasts. The root mean squared errors are 2.32~\%, 1.73~\% and 2.64~\%.
\begin{enumerate}[resume]
\item Which forecast won? Which lost?
\item Find the random walk's and the uncovered-parity forecasts for 2023 and their
errors.
\item Does the result contradict Meese and Rogoff?
\item Why did the PPP rule do badly in 2022?
\item What would you need to say that one forecast is better than another?
\item State the result: the overshoot of the dalen, in points, and the half-life of
its real deviation.
\end{enumerate}
\end{problem}
